\documentclass[a4paper,12pt]{article}

\usepackage[T2A]{fontenc}
\usepackage[utf8]{inputenc}
\usepackage[english,russian]{babel}
\usepackage{geometry}
\usepackage{setspace}
\usepackage{indentfirst}
\usepackage{amsmath}
\usepackage{amssymb}
\usepackage{array}
\usepackage{booktabs}
\usepackage{longtable}
\usepackage{tabularx}
\usepackage{enumitem}
\usepackage{hyperref}
\usepackage{csquotes}
\usepackage{float}
\usepackage[
  backend=bibtex,
  style=numeric,
  sorting=none,
  giveninits=true,
  doi=true,
  url=true,
  isbn=false,
  eprint=false
]{biblatex}

\geometry{left=2.5cm,right=1.8cm,top=2cm,bottom=2cm}
\setlength{\parindent}{1.25cm}
\setstretch{1.35}
\setlength{\LTpre}{0pt}
\setlength{\LTpost}{0pt}
\renewcommand{\arraystretch}{1.15}
\hypersetup{colorlinks=true,linkcolor=black,urlcolor=blue,citecolor=black}
\addbibresource{coursework.bib}

\newcolumntype{L}[1]{>{\raggedright\arraybackslash}p{#1}}

\begin{document}

\begin{titlepage}

{\setstretch{1.0}
\begin{center}
ПРАВИТЕЛЬСТВО РОССИЙСКОЙ ФЕДЕРАЦИИ\\
Федеральное государственное автономное образовательное учреждение высшего образования\\
Национальный исследовательский университет <<Высшая школа экономики>>\\
\bigskip
Факультет компьютерных наук\\
Образовательная программа <<Программная инженерия>>
\end{center}
}

\vspace{4em}

\begin{center}
УДК 004.023

\vspace{2em}

\textbf{Исследование в рамках дисциплины <<Курсовой проект>>}

\vspace{1em}

\textbf{на тему}

\vspace{0.5em}

\textbf{<<Оптимизация алгоритмов маршрутизации на примере задачи нескольких коммивояжеров>>}
\end{center}

\vspace{5em}

\noindent
\begin{tabular}{p{0.42\linewidth}p{0.5\linewidth}}
Научный руководитель & Ермаков Андрей Максимович \\
Исполнитель & студент группы БПИ246 \\
& Д.\,Д. Купцов \\
\end{tabular}

\vspace{\fill}

\begin{center}
Москва -- 2026
\end{center}

\end{titlepage}

\setcounter{page}{2}

{
  \hypersetup{linkcolor=black}
  \tableofcontents
}

\newpage

\section{Введение}

Настоящая работа посвящена исследованию эвристических алгоритмов маршрутизации для задачи множественного коммивояжёра. В проекте рассматривается постановка mTSP с общей вершиной-депо и целевой функцией \texttt{MINSUM}, то есть минимизацией суммарной длины всех маршрутов. Практическая часть исследования ведётся в репозитории \texttt{mtsp-routing-optimization}, где объединены программная реализация, экспериментальный контур и материалы к курсовой работе.

\subsection{Актуальность}

Задачи маршрутизации занимают важное место в комбинаторной оптимизации и напрямую связаны с логистикой, транспортным планированием, доставкой, распределением задач между мобильными агентами и цифровыми сервисами управления инфраструктурой. Классическая задача коммивояжёра (TSP) заключается в поиске кратчайшего замкнутого маршрута, проходящего через все вершины ровно один раз, и давно рассматривается как одна из базовых NP-трудных задач \cite{little1963,junger1995}.

Во многих прикладных сценариях одного маршрута недостаточно: необходимо распределить множество точек между несколькими исполнителями. В этом случае возникает задача множественного коммивояжёра (mTSP), которая по сравнению с классической задачей требует одновременно решать две связанные подзадачи: распределение клиентов между агентами и упорядочение посещения вершин внутри каждого маршрута \cite{bektas2006}. Рост числа вершин и числа агентов быстро увеличивает пространство допустимых решений, поэтому точные методы оказываются трудноприменимыми, а основной практический интерес смещается к эвристикам и метаэвристикам.

Для учебно-исследовательского проекта это особенно важно. С одной стороны, тема позволяет опереться на классические результаты по TSP и методам локального поиска. С другой стороны, задача mTSP даёт содержательную среду для проектирования и сравнения собственных алгоритмов, где качество решения и ограничение по времени вычислений одинаково существенны.

\subsection{Объект и предмет исследования}

Объектом исследования являются процессы маршрутизации и распределения вершин между несколькими агентами в задачах комбинаторной оптимизации.

Предмет исследования составляют эвристические и метаэвристические методы решения задачи множественного коммивояжёра, а также способы их программной реализации и экспериментального сравнения по критерию \texttt{MINSUM}.

\subsection{Методы исследования}

Основным методом исследования является вычислительный эксперимент. В рамках работы используются:
\begin{enumerate}[label=\arabic*.]
  \item анализ научной литературы по задачам TSP и mTSP;
  \item проектирование и реализация эвристических алгоритмов;
  \item сравнительное тестирование на синтетических наборах данных;
  \item анализ качества решений и времени работы алгоритмов;
  \item статистическая интерпретация результатов на одинаковых наборах экземпляров.
\end{enumerate}

Практическая часть исследования строится вокруг воспроизводимого программного контура: генератор инстансов формирует тестовые задачи, единый runner запускает solvers в согласованном формате, а результаты автоматически сохраняются в CSV-артефакты для последующего анализа.

\subsection{Цель исследования}

Цель исследования состоит в разработке и сравнительном анализе эвристических методов оптимизации маршрутов для задачи множественного коммивояжёра при фиксированном ограничении на время вычислений.

\subsection{Задачи исследования}

Для достижения поставленной цели в работе решаются следующие задачи:
\begin{enumerate}[label=\arabic*.]
  \item изучить постановку классической задачи TSP и основные подходы к её решению;
  \item проанализировать особенности постановки mTSP и обзор существующих методов;
  \item исследовать методы локального поиска, GRASP-подход и алгоритмы класса Lin--Kernighan--Helsgaun;
  \item реализовать и привести к единому интерфейсу набор solver'ов для mTSP;
  \item подготовить воспроизводимый экспериментальный контур и набор тестовых инстансов;
  \item провести сравнительные вычислительные эксперименты;
  \item сформулировать выводы о том, какие подходы дают наилучшее качество решения в рамках заданного бюджета времени.
\end{enumerate}

\subsection{Гипотеза исследования}

Рабочая гипотеза проекта заключается в том, что при решении задачи mTSP в условиях ограниченного бюджета времени эвристики, сочетающие построение начального решения и последующее локальное улучшение, дают лучшее значение \texttt{MINSUM}, чем простые baseline-стратегии. В прикладной формулировке это означает, что более сильные методы уровня \texttt{grasp} и \texttt{lkh-wrapper-v2} должны систематически превосходить \texttt{rand+nn} и \texttt{2opt+greed} по качеству маршрутов, а их преимущество должно сохраняться при росте размерности задачи и числа коммивояжёров.

\section{Новизна и достоверность результатов}

Новизна работы состоит в том, что задача mTSP рассматривается не только как объект литературного обзора, но и как основа для построения собственного воспроизводимого экспериментального контура. В проекте объединяются несколько уровней исследования: единая программная архитектура solver'ов, синтетический benchmark-набор, сравнение с внешним reference baseline и оформление результатов в виде связанного корпуса материалов к курсовой работе.

Важно отметить, что в текущей версии проекта используется собственная эвристика \texttt{lkh-wrapper-v2}, вдохновлённая идеями Lin--Kernighan--Helsgaun, а не прямое подключение внешнего пакета LKH-3. Такое уточнение делает описание работы методически корректным: исследование посвящено сравнению реально реализованных алгоритмов, а не декларации ещё не интегрированных компонентов.

Достоверность результатов обеспечивается несколькими обстоятельствами. Во-первых, все solver'ы запускаются на одном и том же наборе инстансов и через единый программный интерфейс. Во-вторых, результаты автоматически фиксируются в машинно-читаемых CSV-файлах, что снижает риск ручных ошибок и позволяет проверять расчёты повторно. В-третьих, для интерпретации качества решений используются одинаковые метрики и сопоставление на парных наборах экземпляров, что соответствует общим рекомендациям по экспериментальному анализу алгоритмов \cite{hooker1995,johnson1999}.

\subsection{Теоретическая значимость результатов}

Теоретическая значимость работы состоит в систематизации подходов к решению задачи mTSP и в аккуратном разведении трёх уровней анализа:
\begin{enumerate}[label=\arabic*.]
  \item постановка и свойства задачи;
  \item структура и ограничения конкретных эвристических алгоритмов;
  \item экспериментальный дизайн и статистическая интерпретация результатов.
\end{enumerate}

Такое разделение важно для курсовой работы, поскольку позволяет не смешивать корректность реализованных процедур с выводами об их сравнительной эффективности.

\subsection{Практическая ценность результатов}

Практическая ценность работы определяется тем, что её результатом является не только текстовое описание, но и рабочая программная система. В репозитории уже реализованы базовые и усиленные solver'ы, генерация синтетических инстансов, сценарии массовых прогонов, формирование сводных отчётов и отдельные приложения к курсовой. Это означает, что результаты могут использоваться как для дальнейшего развития проекта, так и для подготовки финальной версии исследования без пересборки всего контура с нуля.

\section{Обзор и анализ источников, выбор методов для решения поставленных задач}

\subsection{Обзор существующих подходов к решению задачи TSP}

Задача коммивояжёра является одной из наиболее известных задач дискретной оптимизации. Для неё разработаны как точные методы, так и многочисленные эвристические схемы \cite{little1963,junger1995}. На небольших размерностях применимы алгоритмы ветвей и границ, динамического программирования и целочисленного программирования, однако их стоимость быстро растёт вместе с числом вершин. Поэтому в прикладных исследованиях широко используются конструктивные эвристики, локальный поиск и гибридные схемы улучшения.

Особое место среди классических методов занимает семейство \(k\)-opt-перестановок. Для практики наиболее важен 2-opt, в котором удаляются два ребра маршрута и выполняется переподключение с уменьшением длины цикла \cite{croes1958}. Этот подход не гарантирует глобального оптимума, но даёт простой и эффективный механизм локального улучшения.

\subsection{Подходы к решению задачи множественного коммивояжёра mTSP}

Задача mTSP является естественным обобщением TSP на случай нескольких агентов. В отличие от классической постановки, здесь важно не только определить порядок обхода вершин, но и распределить вершины между маршрутами \cite{bektas2006,russell1977}. В зависимости от постановки могут минимизироваться суммарная длина маршрутов, максимальная длина одного маршрута или иные критерии баланса. В данной работе рассматривается вариант \texttt{MINSUM}, в котором минимизируется суммарная длина всех маршрутов.

Для решения mTSP применяются жадные схемы, кластеризационные подходы, локальный поиск, имитация отжига, табу-поиск и гибридные метаэвристики. Общий вывод литературного обзора заключается в том, что для задач средней и большой размерности основную роль играют именно эвристические методы, позволяющие получать качественные решения в разумное время.

\subsection{Алгоритмы класса Lin--Kernighan--Helsgaun и их применение к работе}

Алгоритм Lin--Kernighan является одним из важнейших результатов в области TSP: он использует адаптивную стратегию выбора перестроек вместо фиксированного значения \(k\) в \(k\)-opt-поиске \cite{lin1973}. Впоследствии эти идеи были существенно усилены в работе К. Хельсгауна, где предложены кандидатные множества, более эффективные правила выбора ходов и развитый механизм локального поиска \cite{helsgaun2000}.

Для настоящего проекта важен не столько буквальный вызов внешнего LKH-пакета, сколько сами принципы построения сильной эвристики: предварительный отбор перспективных соседей, содержательная конструктивная фаза, повторное локальное улучшение и межмаршрутная дооптимизация. Именно по этой логике в репозитории развивалась собственная линия \texttt{lkh-wrapper}.

\subsection{Выбор методов для решения поставленных задач}

С учётом литературного обзора и текущего состояния репозитория в работе зафиксированы четыре основных solver'а:
\begin{enumerate}[label=\arabic*.]
  \item \texttt{rand+nn} как нижняя baseline-планка: случайное распределение клиентов между агентами с последующим nearest-neighbor построением;
  \item \texttt{2opt+greed} как базовая схема <<конструирование + локальное улучшение>>;
  \item \texttt{grasp} как стохастический метаэвристический метод с жадно-рандомизированным построением и локальным поиском \cite{feo1995};
  \item \texttt{lkh-wrapper-v2} как наиболее сильная текущая эвристика проекта, сочетающая candidate-based выбор, lookahead-оценку, 2-opt и межмаршрутные обмены.
\end{enumerate}

Такой набор методов позволяет решать сразу две исследовательские задачи: сравнить слабые и сильные эвристики внутри одной кодовой базы и оценить, насколько собственные решения приближаются к внешнему reference baseline на базе OR-Tools.

\section{Журнал эксперимента}

Экспериментальная часть проекта оформлена как воспроизводимый вычислительный контур. Его задача состоит не только в разовом запуске solver'ов, но и в накоплении сравнимых артефактов, по которым можно последовательно развивать курсовую работу. На текущем этапе часть выводов уже подтверждена основными CSV-файлами, а часть правил зафиксирована как итоговый протокол финальной версии исследования.

\subsection{Режимы эксперимента}

\begin{table}[H]
\centering
\caption{Основные режимы экспериментального контура}
\begin{tabular}{L{0.22\linewidth}L{0.27\linewidth}L{0.41\linewidth}}
\toprule
Режим & Назначение & Состав \\
\midrule
Основной текущий benchmark & Сравнение реализованных solver'ов на единой сетке экземпляров & \texttt{uniform}-инстансы, \(n \in \{20,50,100,200\}\), \(m \in \{2,3,5\}\), \(10\) экземпляров на комбинацию параметров \\
Pilot по повторным seed & Оценка разброса стохастических методов и выбор числа повторов & Выбранные инстансы семейств \texttt{uniform}, \texttt{clustered-center}, \texttt{clustered-offset-depot}, \texttt{mixed-outliers}; независимые повторные запуски \\
Reference baseline & Сопоставление с внешним сильным ориентиром по качеству решения & OR-Tools в режиме \texttt{GUIDED\_LOCAL\_SEARCH} с отдельным CSV артефактом \\
Целевой финальный benchmark & Итоговая схема для полной версии курсовой & Четыре семейства инстансов, базовый объём \(N=10\) на страту, подтверждающий режим \(N=20\), дополнительный stress-режим \(m=7\) для \(n=100\) и \(n=200\) \\
\bottomrule
\end{tabular}
\end{table}

Текущая конфигурация генератора уже поддерживает четыре семейства синтетических инстансов: \texttt{uniform}, \texttt{clustered-center}, \texttt{clustered-offset-depot} и \texttt{mixed-outliers}. Однако подтверждённый основной CSV на момент подготовки данного текста описывает прежде всего \texttt{uniform}-сетку. Поэтому в текущей редакции раздела важно различать уже полученные результаты и тот полный benchmark, который закреплён как финальный протокол дальнейшей работы.

\subsection{Последовательность выполнения}

В проекте используется следующая последовательность вычислительного эксперимента:
\begin{enumerate}[label=\arabic*.]
  \item генерация набора mTSP-инстансов с помощью скрипта \path{experiments/generate_mtsp_instances.py};
  \item запуск solver'ов через единый runner \path{run_mtsp.py} и пакетный сценарий \path{experiments/run_benchmarks.py};
  \item автоматическая запись сырых результатов в \path{data/results/mtsp_results.csv};
  \item построение агрегированной сводки \path{data/results/mtsp_summary.csv};
  \item отдельный расчёт comparison with external reference baseline и сохранение результатов в \path{data/results/mtsp_reference_results.csv} и \path{data/results/mtsp_reference_summary.csv};
  \item парный анализ различий между solver'ами и pilot-оценка разброса для стохастических методов.
\end{enumerate}

Такая схема позволяет отделить уровень <<получить решение на одном экземпляре>> от уровня <<собрать сопоставимый набор результатов по всему benchmark-набору>>.

\subsection{Повторы и выбор итогов}

Для детерминированных solver'ов по качеству решения достаточно одного запуска на фиксированном экземпляре. Для стохастических методов этого недостаточно, поэтому в итоговом протоколе работы закреплены независимые повторные запуски по разным значениям \texttt{seed}. На текущем этапе основной CSV ещё строился с фиксированным \texttt{seed=42} для стохастических solver'ов, тогда как pilot уже использует серию независимых повторов и позволяет оценить разброс результатов.

В качестве базового правила для дальнейшего расширения исследования принимается следующая схема:
\begin{enumerate}[label=\arabic*.]
  \item основной режим: \(N=10\) экземпляров на страту и \(R=20\) повторов для стохастических solver'ов;
  \item подтверждающий режим: \(N=20\) и \(R=30\) для спорных сравнений;
  \item строгий локальный режим: \(R=50\), если требуется уточнить результат в отдельных комбинациях параметров.
\end{enumerate}

Итоговой оценкой качества стохастического solver'а на фиксированном экземпляре считается среднее по независимым повторам. Для попарного сравнения алгоритмов используются одни и те же экземпляры внутри заданной страты, что делает выводы более устойчивыми и статистически корректными.

\subsection{Агрегирование результатов}

В работе используются следующие уровни агрегирования:
\begin{enumerate}[label=\arabic*.]
  \item метрики одного запуска: \texttt{objective}, \texttt{time}, \texttt{valid}, набор маршрутов и история шагов;
  \item сводка по solver'у и страте: среднее значение \texttt{MINSUM}, среднее время работы, число валидных прогонов;
  \item сравнение с reference baseline: абсолютный разрыв по целевой функции и относительный gap в процентах;
  \item парные разности между solver'ами на одном и том же наборе экземпляров.
\end{enumerate}

Основной смысл такого агрегирования состоит в том, чтобы отдельно анализировать качество решения, вычислительную стоимость и устойчивость результатов, а не сводить их к одной неинтерпретируемой оценке.

\subsection{Фиксация артефактов}

Ключевые артефакты проекта, на которые опирается текущая версия курсовой, перечислены ниже:
\begin{enumerate}[label=\arabic*.]
  \item конфигурация основного benchmark-а: \path{experiments/config.json};
  \item сырые результаты прогонов: \path{data/results/mtsp_results.csv};
  \item агрегированная сводка: \path{data/results/mtsp_summary.csv};
  \item результаты сравнения с OR-Tools: \path{data/results/mtsp_reference_results.csv}, \path{data/results/mtsp_reference_summary.csv};
  \item текущая парная статистика: \path{data/results/methodology_pairwise_current.csv};
  \item pilot по повторным запускам: \path{data/results/methodology_pilot_summary.csv}.
\end{enumerate}

В окончательную редакцию раздела будут добавлены таблицы по полному четырёхсемейному benchmark-набору после завершения соответствующих прогонов.

\section{Выводы эксперимента}

На текущем подтверждённом наборе результатов рабочая гипотеза получила содержательную поддержку. По данным сводного отчёта сравнения с внешним reference baseline средний относительный gap к OR-Tools минимален у \texttt{lkh-wrapper-v2}. За ним следуют \texttt{grasp}, \texttt{2opt+greed} и \texttt{rand+nn}. Это означает, что усиленные эвристики действительно дают более качественные решения, чем простые baseline-методы.

\begin{table}[H]
\centering
\caption{Текущая агрегированная сводка по uniform-набору}
\begin{tabular}{L{0.25\linewidth}L{0.17\linewidth}L{0.17\linewidth}L{0.18\linewidth}}
\toprule
Solver & Средний gap, \% & Средний абсолютный gap & Среднее время, c \\
\midrule
\texttt{lkh-wrapper-v2} & 14.48 & 92.60 & 0.0181 \\
\texttt{grasp} & 20.77 & 138.19 & 1.0469 \\
\texttt{2opt+greed} & 47.41 & 314.21 & 0.0007 \\
\texttt{rand+nn} & 114.73 & 843.01 & 0.0005 \\
\bottomrule
\end{tabular}
\end{table}

Полученные результаты важны не только на уровне общей средней. По текущей uniform-сетке \texttt{lkh-wrapper-v2} показывает лучший средний gap в \(10\) из \(12\) страт, а \texttt{grasp} лидирует лишь в двух сравнительно простых комбинациях параметров. Парная статистика между \texttt{grasp} и \texttt{lkh-wrapper-v2} также подтверждает преимущество более сильного solver'а: в \(10\) из \(12\) страт средняя разность благоприятствует \texttt{lkh-wrapper-v2}, а в \(8\) из \(12\) страт \(95\%\)-й доверительный интервал для этой разности уже не содержит нуль.

Отдельный pilot по повторным запускам показывает, что разброс результатов у \texttt{lkh-wrapper-v2} существенно ниже, чем у слабых baseline-подходов, а для \texttt{grasp} разумный контроль дисперсии достигается при \(R=20\) или \(R=30\) независимых повторах. Следовательно, итоговый протокол с повторными seed-запусками не является формальностью: он нужен для корректного сравнения стохастических solver'ов и для окончательной фиксации статистически устойчивых выводов.

При этом результаты текущего раздела нужно интерпретировать аккуратно. Основной подтверждённый CSV пока отражает прежде всего \texttt{uniform}-инстансы и фиксированный \texttt{seed=42} в базовом прогоне. Поэтому финальные выводы всей курсовой должны опираться на полный benchmark-набор по четырём семействам данных и на расширенный режим повторных запусков. Тем не менее уже сейчас можно сделать промежуточный вывод: переход от простых baseline-эвристик к более сильным методам действительно улучшает качество решения задачи mTSP.

\section{Техническое описание программной системы эксперимента}

\subsection{Назначение системы}

Программная система проекта предназначена для воспроизводимого исследования эвристических алгоритмов решения задачи множественного коммивояжёра. Она должна обеспечивать три функции: запуск solver'ов в едином формате, сбор и сохранение результатов, а также подготовку материалов для анализа и оформления курсовой работы.

\subsection{Основные функции системы}

Система выполняет следующие основные функции:
\begin{enumerate}[label=\arabic*.]
  \item загрузка mTSP-инстансов из текстового формата;
  \item запуск выбранного solver'а или цепочки solver'ов через единый CLI-интерфейс;
  \item проверка валидности построенных маршрутов;
  \item вычисление значения целевой функции и времени работы;
  \item пакетный прогон алгоритмов на benchmark-наборе;
  \item формирование CSV-артефактов и сводных таблиц для дальнейшего анализа.
\end{enumerate}

\subsection{Архитектура и состав модулей}

\begin{table}[H]
\centering
\caption{Ключевые компоненты программной системы}
\begin{tabular}{L{0.29\linewidth}L{0.53\linewidth}}
\toprule
Компонент & Назначение \\
\midrule
\path{src/mtsp_main.cpp} & C++-точка входа для запуска mTSP solver'ов и формирования JSON-результата \\
\path{src/mtsp_rand_nn_solver.cpp} & baseline-эвристика случайного распределения с nearest-neighbor построением \\
\path{src/mtsp_greedy_2opt_solver.cpp} & жадная конструктивная схема с последующим 2-opt-улучшением \\
\path{src/mtsp_grasp_solver.cpp} & стохастическая реализация GRASP-подхода \\
\path{src/mtsp_lkh_wrapper_v2.cpp} & усиленная эвристика проекта, использующая candidate-based выбор и локальный поиск \\
\path{python/mtsp_runner.py} и \path{run_mtsp.py} & Python-обёртка для запуска C++-решателя на одном экземпляре \\
\path{experiments/generate_mtsp_instances.py} & генерация синтетических benchmark-инстансов \\
\path{experiments/run_benchmarks.py} & пакетный прогон solver'ов и запись результатов в CSV \\
\path{data/results/} & каталог итоговых численных артефактов и сводок \\
\bottomrule
\end{tabular}
\end{table}

Архитектурно проект устроен как связка <<C++-ядро + Python-оркестрация + LaTeX-документация>>. Вычислительно тяжёлая часть вынесена в C++-модули solver'ов, а сценарии массовых запусков, агрегации и подготовки отчётных артефактов реализованы на Python. Такой подход упрощает расширение экспериментального контура и делает систему удобной для последующего сопровождения.

\subsection{Входные и выходные данные}

Базовый формат входного mTSP-инстанса имеет следующий вид:

\begin{verbatim}
n m
x0 y0
x1 y1
...
\end{verbatim}

Здесь \(n\) обозначает число вершин, \(m\) --- число коммивояжёров, а вершина с индексом \(0\) рассматривается как общее депо. Такой формат прост для генерации, чтения и последующего пакетного использования в экспериментах.

На выходе solver формирует JSON-объект, содержащий:
\begin{enumerate}[label=\arabic*.]
  \item маршруты;
  \item итоговое значение целевой функции;
  \item общее время выполнения;
  \item признак валидности решения;
  \item при необходимости историю промежуточных шагов.
\end{enumerate}

При пакетных прогонах эти данные автоматически преобразуются в CSV-таблицы, пригодные для статистической обработки и включения в текст курсовой работы.

\subsection{Сценарий использования в эксперименте}

Типовой сценарий использования системы в рамках исследования выглядит следующим образом:
\begin{enumerate}[label=\arabic*.]
  \item формируется набор синтетических инстансов заданных семейств и размерностей;
  \item для каждого экземпляра запускаются все solver'ы из конфигурации эксперимента;
  \item результаты каждого прогона сохраняются в сырую таблицу;
  \item поверх сырых данных строятся агрегированные сводки и сравнение с внешним reference baseline;
  \item по полученным артефактам формируются текстовые выводы и таблицы для курсовой.
\end{enumerate}

Таким образом, программная система служит не только инструментом вычислений, но и основой исследовательской дисциплины проекта: один и тот же эксперимент можно повторить, проверить и дополнить без изменения общей логики работы.

\section{Основные термины и сокращения}

\begin{longtable}{L{0.27\linewidth}L{0.63\linewidth}}
\toprule
Термин & Определение \\
\midrule
\endfirsthead
\toprule
Термин & Определение \\
\midrule
\endhead
\texttt{TSP} & Задача коммивояжёра, в которой требуется построить один кратчайший замкнутый маршрут через все вершины. \\
\texttt{mTSP} & Задача множественного коммивояжёра, в которой множество вершин распределяется между несколькими маршрутами с общим депо. \\
\texttt{MINSUM} & Целевая функция, равная суммарной длине всех маршрутов; основной критерий качества в данной работе. \\
Solver & Реализация конкретного алгоритма решения задачи маршрутизации в составе программной системы. \\
Baseline-метод & Простой опорный алгоритм, относительно которого оценивается выигрыш более сильных эвристик. \\
Локальный поиск & Последовательное улучшение текущего решения через переходы к соседним решениям. \\
\texttt{2-opt} & Частный случай \(k\)-opt-перестановки, при котором перестраиваются два ребра маршрута для сокращения его длины. \\
\texttt{GRASP} & Greedy Randomized Adaptive Search Procedure, метаэвристический подход, сочетающий жадное рандомизированное построение и локальное улучшение \cite{feo1995}. \\
\texttt{LKH} & Класс сильных эвристик для TSP, основанных на идеях Lin--Kernighan и развитых в работах Хельсгауна \cite{lin1973,helsgaun2000}. \\
Reference baseline & Внешний ориентир для сравнения качества решений; в проекте эту роль выполняет OR-Tools, но не как доказанный оптимум. \\
Синтетический инстанс & Искусственно сгенерированная задача, используемая для контролируемого сравнения алгоритмов. \\
Страта эксперимента & Комбинация параметров семейства инстансов, числа вершин и числа коммивояжёров, внутри которой сравниваются solver'ы. \\
\bottomrule
\end{longtable}

\section{Календарный план работ}

\begin{longtable}{L{0.2\linewidth}L{0.35\linewidth}L{0.22\linewidth}L{0.15\linewidth}}
\toprule
Этап проекта & Содержание работ & Ожидаемый результат & Сроки \\
\midrule
\endfirsthead
\toprule
Этап проекта & Содержание работ & Ожидаемый результат & Сроки \\
\midrule
\endhead
Анализ TSP и базовых эвристик & Изучение постановки классической задачи коммивояжёра, обзор точных и эвристических методов, анализ жадных схем и локального поиска. & Теоретическая база для дальнейшего исследования. & 06.11.2025 -- 30.11.2025 \\
Изучение mTSP и литературного фона & Анализ постановки mTSP, обзор существующих эвристических подходов, выбор целевой функции \texttt{MINSUM}. & Формулировка предмета исследования и рабочей гипотезы. & 01.12.2025 -- 31.12.2025 \\
Проектирование программного контура & Выбор архитектуры проекта, унификация формата входа и выхода, подготовка общего runner'а для mTSP solver'ов. & Базовая программная инфраструктура эксперимента. & 01.01.2026 -- 31.01.2026 \\
Подготовка данных и benchmark-набора & Реализация генератора синтетических инстансов, настройка конфигурации эксперимента, фиксация основного набора solver'ов. & Воспроизводимый benchmark-набор и сценарии запусков. & 01.02.2026 -- 15.03.2026 \\
Проведение вычислительных экспериментов & Массовые прогоны solver'ов, сбор результатов, сравнение с reference baseline, pilot по повторным seed-запускам. & Набор CSV-артефактов и предварительные выводы по качеству алгоритмов. & 16.03.2026 -- 10.04.2026 \\
Оформление курсовой и приложений & Подготовка основного текста курсовой, приложений по математическому обоснованию и методике экспериментов, уточнение итоговых таблиц. & Финальная связная версия исследовательской работы. & 11.04.2026 -- 15.05.2026 \\
\bottomrule
\end{longtable}

\printbibliography[title={Список литературы}]

\end{document}
